% =================================================================
\section{Proposed System}
\label{sec:method}
% =================================================================

Fig.~\ref{fig:system} shows the full pipeline.
We use the MTRCNN backbone~\cite{mtrcnn2025} throughout: three parallel
convolutional branches (temporal kernel sizes 3, 5, 7 capture different
wingbeat timescales), each with three Conv-BN-ReLU-Pool stages,
followed by masked mean$+$max pooling that ignores padding frames.
The three branch outputs are concatenated into a 32-dimensional
shared embedding~$z$ ($\approx$221K parameters total; the 32-dim size
follows the original MTRCNN paper).
Two linear heads predict species (9 classes) and recording domain
(5 classes) from~$z$.

\subsection{Domain-Balanced Sampling}
\label{sec:method:balance}

\textbf{The problem it solves:} without intervention, D5 appears in
99.4\% of every training batch, so the model learns almost exclusively
from one recording condition.
The other four conditions are effectively invisible.

\textbf{The fix:} we assign each training clip a sampling weight
$w_i = 1 / \text{count}(d_i)$, where $d_i$ is its domain label.
This equalises how often each domain is seen per epoch, upsampling
D1--D4 clips by 200--2{,}600$\times$.
As discussed in Section~\ref{sec:data:imbalance}, every subsequent
method in this paper depends on this step: without balanced training,
the model has no D1--D4 signal to generalise from.

\subsection{Domain-Adversarial Training (DANN)}
\label{sec:method:dann}

\textbf{The problem it solves:} even with balanced training, the
embedding~$z$ may still encode recording-condition information
alongside species information, letting the model ``cheat'' by
memorising condition-specific cues.

\textbf{How it works:} a Gradient Reversal Layer~(GRL)~\cite{dann2016}
sits between~$z$ and the domain head.
During the forward pass it is an identity; during backpropagation it
flips the sign and scales the gradient by~$\lambda(p)$.
This makes the embedding \emph{worse} at predicting the domain---the
model is forced to discard recording-condition information to keep the
domain head from working.
The schedule from~\cite{dann2016} gradually increases the adversarial
pressure:
\begin{equation}
  \lambda(p) = \lambda_{\max}\!\cdot\!\left(\frac{2}{1+e^{-10p}}-1\right),
  \quad p = \frac{e}{E},
  \label{eq:dann_schedule}
\end{equation}
where $e$ is the current epoch and $E{=}100$ (max epochs).
The constant~10 controls how quickly the adversary ramps up~\cite{dann2016};
at the halfway point ($p{=}0.5$) the model is already at 96\% of its
full adversarial pressure.
We set $\lambda_{\max}{=}1.0$.

\subsection{Domain-Invariant Contrastive Learning (DiCL)}
\label{sec:method:dicl}

\textbf{The problem it solves:} DANN pushes domain information
\emph{out} of the embedding, but it does not explicitly pull
same-species recordings from different conditions
\emph{together}.

\textbf{How it works:} DiCL~\cite{drbiol2025} treats any two clips of
the same species from \emph{different} recording conditions as a
positive pair.
For each such anchor~$i$, the contrastive loss pulls its embedding
toward all its cross-condition same-species partners, while pushing
it away from all other clips in the batch.

Concretely, let $P_i = \{j : y_j = y_i,\, d_j \neq d_i\}$ be the
cross-condition positives for anchor~$i$, and let
$\mathcal{A} = \{i : |P_i| > 0\}$.
Embeddings are first passed through a two-layer projection head
($h_i = \mathrm{proj}(z_i) \in \mathbb{R}^{128}$; architecture:
$32{\to}128{\xrightarrow{\mathrm{ReLU}}}128$, no final activation).
The projection head means the contrastive geometry is learned in a
separate space from the classifier, which prevents the loss from
distorting the species-discriminative directions in~$z$~\cite{supcon2020}.
DiCL then minimises:
\begin{equation}
  \mathcal{L}_{\mathrm{DiCL}} =
  {-}\frac{1}{|\mathcal{A}|}\!\sum_{i \in \mathcal{A}}
  \frac{1}{|P_i|}\!\sum_{j \in P_i}
  \log\frac{e^{\,\tilde{h}_i \cdot \tilde{h}_j / \tau}}
           {\sum_{k \neq i} e^{\,\tilde{h}_i \cdot \tilde{h}_k / \tau}},
  \label{eq:dicl}
\end{equation}
where $\tilde{h} = h / \|h\|_2$.
We use temperature $\tau{=}0.2$ (rather than 0.07 in~\cite{drbiol2025}):
with only 80--634 minority-domain clips per epoch there are very few
positive pairs per batch, and a lower temperature concentrates the
gradient on only the hardest negatives, making learning unstable.
The combined training loss is:
\begin{equation}
  \mathcal{L} = \mathcal{L}_{\mathrm{sp}} - \lambda(p)\,\mathcal{L}_{\mathrm{dom}}
              + w_c\,\mathcal{L}_{\mathrm{DiCL}},
  \label{eq:total_loss}
\end{equation}
with $w_c{=}0.1$~\cite{drbiol2025} and cross-entropy losses
$\mathcal{L}_{\mathrm{sp}}$, $\mathcal{L}_{\mathrm{dom}}$ for species
and domain.

\subsection{Frequency-Band-Selective Mix (FBS-Mix)}
\label{sec:method:fbsmix}

\textbf{The problem it solves:}
MixStyle~\cite{mixstyle2021} diversifies domain characteristics by
randomly blending the mean and variance of CNN feature maps between
training samples.
It works well on vision tasks, but here it backfires: it mixes the
frequency bands that carry \emph{species} information (the wingbeat
tone) at the same time as the recording-noise bands it is trying to
diversify.
The two effects cancel, and balanced+MixStyle performs no better than
balanced alone ($\approx$0.32~$\BAu$).

\textbf{The insight:}
the per-frequency variance analysis~(Section~\ref{sec:data:freq_analysis})
tells us exactly which bands to mix and which to leave alone.
Bins 0--8 ($\approx$36--325~Hz) are dominated by recording-condition
noise.
Bins 9--35 ($\approx$361--1{,}360~Hz) carry the mosquito wingbeat and
are species-discriminative.
So we mix \emph{only} bins 0--8 and protect the rest.

\textbf{How it works:}
FBS-Mix is an input-level batch transform, applied before the model
with probability $p{=}0.5$ per batch (stochastic application
prevents the model from depending on the augmentation).
It implements AdaIN~\cite{huang2017adain}---adaptive instance
normalisation---on the domain-noise band only.
Given a padded batch $X \in \mathbb{R}^{B \times T \times F}$, extract
the low-frequency band $X_{\mathrm{lo}} = X_{:,:,0:K}$ ($K{=}9$).
Sample a mixing coefficient $\lambda \sim \mathrm{Beta}(\alpha,\alpha)$
with $\alpha{=}0.1$ (a U-shaped distribution: usually $\lambda$ is near
0 or 1, so most batches see a clear style donor, and blending is
occasionally subtle) and a random batch permutation~$\pi$.
Compute per-sample, per-bin statistics $(\mu_i,\sigma_i)$ over
\emph{valid frames only}---padding frames are excluded, because a
63-frame clip padded to 200 frames is 68\% zeros and would severely
bias na\"ive global statistics.
Transfer the blended statistics to the band:
\begin{equation}
  \tilde{X}_{\mathrm{lo},i} =
    \underbrace{\bigl(\lambda\,\sigma_i + (1{-}\lambda)\,\sigma_{\pi(i)}\bigr)}_{\text{blended std}}
    \cdot \frac{X_{\mathrm{lo},i} - \mu_i}{\sigma_i}
    +\; \underbrace{\bigl(\lambda\,\mu_i + (1{-}\lambda)\,\mu_{\pi(i)}\bigr)}_{\text{blended mean}}.
  \label{eq:fbsmix}
\end{equation}
Bins $K{:}K'$ (wingbeat, $K'{=}36$, 361--1{,}360~Hz) and bins
$K'{:}F$ (upper harmonics and noise, 1{,}412--3{,}854~Hz) are left
unchanged.
Because FBS-Mix operates on the raw spectrogram before any model
layers, there is no inference-time overhead: it is simply switched
off during evaluation.
